\documentclass[ 11pt
				,ngerman 
				,headsepline
				,headings=small
				,numbers=noenddot % Kein "." hinter der letzten Kapitelnummer --> 2.1 statt 2.1.
				,draft=false
				,BCOR=0mm % Wert für die Bindung 
				,DIV=12 % Wert für die größe des Textblocks
				,captions=tableheading
				,paper=a4
				,abstracton
				,twoside % Erlaubt beidseitigen Druck
                ]{scrreprt}
% Wählt den passenden Offset und Margin für das Deckblatt und wählt ein passenden Seitenlayout
\makeatletter
\newlength\offsetCoverPage
\newlength\marginCoverPage
\newcount\onesideprint
\if@twoside
\usepackage[a4paper, left=2.6cm, right=2.9cm, top=3.7cm, bottom=3.8cm]{geometry}
\offsetCoverPage=0.0cm
\marginCoverPage=0.25cm
\onesideprint=0
\else
\usepackage[a4paper, left=2.75cm, right=2.75cm, top=3.7cm, bottom=3.8cm]{geometry}
\offsetCoverPage=-0.10cm
\marginCoverPage=0cm
\onesideprint=1
\fi
\makeatother


%-----------------------------------------------------------------------------
% Größeninformationen:
% mit den gewählten Einstellungen beträgt die Textbreite (und damit auch die
% sinnvolle Bildbreite) ca. 158 mm
% Die Texthöhe beträgt etwa 215.0 mm, eine vollseitige Grafik sollte daher 
% maximal 200 mm hoch sein, halbseitige Grafiken entsprechend etwa 90 mm um
% genügend Raum für ein zweites Bild mit Bildunterschrift oder alternativ 
% Fließtext zu bieten
%-----------------------------------------------------------------------------

%-----------------------------------------------------------------------------
% Basispackete
%-----------------------------------------------------------------------------
\usepackage[utf8]{inputenc} 
%\usepackage[english]{babel} %Für einen englischen Text kommentiere die nächste Zeile und entferne den Kommentar dieser Zeile
\usepackage[ngerman]{babel} 
\usepackage{ifpdf} 
\usepackage{blindtext} % Um den lorem ipsum Text zu erzeugen
                      
%-----------------------------------------------------------------------------
% Schrifteinstellungen
%-----------------------------------------------------------------------------
\usepackage{fourier} %PDF-compatible font "Utopia". There exists a free look-alike: Heuristica
\makeatletter % No warnings because of a replaced font
  \let\@font@info\@gobble
  \let\@font@warning\@gobble
\makeatother
\usepackage[official]{eurosym} % Eurosymbol
\DeclareUnicodeCharacter{20AC}{\euro} % The normale euro symbol via [Alt Gr][€] is replaced by a high quality one
\renewcommand{\caplabelfont}{\bfseries} % Fig. and Tab. are bold
\usepackage{color}
\definecolor{ercred}{RGB}{229,48,39}
\usepackage{setspace} % Change Linespread
\onehalfspacing % Linespread 1.5
\setlength\parskip{\medskipamount} % Small linespread between two paragraphs
\setlength\parindent{0pt} % No indentation in the first line of a new paragraph
\AtBeginDocument{
  \renewcommand{\labelitemi}{\(\triangleright\)} % re-definition of the bullet-points
  \renewcommand{\labelitemii}{\(\bullet\)}}

%-----------------------------------------------------------------------------
% Seitenlayout
%-----------------------------------------------------------------------------
\usepackage[automark, autooneside=false]{scrlayer-scrpage}
% Für eine genauere Erklärung lies scrguide.pdf, ca. Seite 240
\automark[chapter]{chapter}
\automark*[section]{}
\interfootnotelinepenalty=10000 % Fußnoten nicht über zwei Seiten verteilen

%Kopfzeile einseitiger Druck
\ifnum\onesideprint=1
\pagestyle{scrheadings}
 \clearscrheadfoot
 \automark{chapter}
 \renewcommand\sectionmark[1]{\markright{\MakeMarkcase {\thesection\hskip .5em\relax#1}}} 
 \rohead{\ifnum\expandafter\pdfstrcmp\botmark=0 \rightmark\else\leftmark{} --- \rightmark\fi}
\ihead{\leftmark}
\ohead{\rightmark}
\cfoot{\pagemark}
\fi
%-----------------------------------------------------------------------------
% Tabellen
%-----------------------------------------------------------------------------
\usepackage{longtable}
\usepackage{multirow}
\usepackage{booktabs} % Erstelle besser aussehende Tabellen, nutze \toprule, \midrule, \bottomrule für horizontale Linien
\usepackage[point]{rccol} % Werte werden über das Kommata ausgerichtet

%-----------------------------------------------------------------------------
% Abbildungen
%-----------------------------------------------------------------------------
\usepackage{subfigure}
\usepackage{float}
% Reduzierung der Anzahl von "Float-Only"-Seiten
\renewcommand{\topfraction}{0.85}
\renewcommand{\textfraction}{0.1}
\renewcommand{\floatpagefraction}{0.75}

%-----------------------------------------------------------------------------
% Mathematik
%-----------------------------------------------------------------------------
\usepackage{amsmath}
\usepackage{amssymb}
\usepackage{array}
\usepackage[cdot,thickqspace,squaren,textstyle]{SIunits} % SI-Einheiten können mit Namen eingegeben werden

%-----------------------------------------------------------------------------
% Verweise und Zitate
%-----------------------------------------------------------------------------
\usepackage{varioref}
\usepackage[square]{natbib}
\usepackage{bibunits}

%-----------------------------------------------------------------------------
% PDF-Setup
%-----------------------------------------------------------------------------
\ifpdf
	%wir benutzen PDFLaTeX
	\usepackage[pdftex]{graphicx} 
	\pdfimageresolution=100
	\pdfminorversion=7
	\pdfcompresslevel=9
	\usepackage[pdftex,
		urlcolor=blue,
		linktoc=all,
		colorlinks=true,
		linkcolor=black,
		citecolor=black,
		pdfview=FitH,
		pdfstartview=FitH,
		plainpages=false]{hyperref}
	\hypersetup{
		pdfauthor={Your Name here},
        pdftitle={Name of your thesis},
        pdfsubject={},
        pdfkeywords={several helpful keywords},
        pdfproducer={LaTeX with hyperref},
        pdfcreator={pdflatex}}
    \usepackage[figure,figure*]{hypcap}
\else
	%DVI oder PS Ausgabe
	\usepackage{graphicx} 
\fi

%-----------------------------------------------------------------------------
% JETZT GEHTS LOS!
%-----------------------------------------------------------------------------
\begin{document}
\pagestyle{empty}
\newcounter{Hilfszaehler}

%-----------------------------------------------------------------------------
% Deckblatt
%-----------------------------------------------------------------------------
\begin{titlepage}
\begin{addmargin}{\offsetCoverPage} 
%TODO Um das RWTH-Logo zu benutzen, musst du das Formblatt RWTH-Logo Deutsch aus der Datei Formulare.txt
%TODO ausfüllen und es dem ZPA schicken!

%-----------------------------------------------------------------------------
% Platzierung der Kopfzeile(Nummer und Logo)
%-----------------------------------------------------------------------------
\setlength{\unitlength}{1cm}
\begin{picture}(0,0)
\put(7.977,0.68){\includegraphics[width = 0.5\textwidth]{Pictures/rwth_eerc_rgb_ohne_Schutzraum}}
%TODO Ändere die Nummer deiner Arbeit, und wenn noetig die Bezeichnung(BA, MA, PA)
%TODO Die Abkürzungen sind Deutsch, auch wenn die Arbeit in Englisch verfasst wird
\put(0,1.32){\fontfamily{phv}\selectfont\huge{BA tba}}
\end{picture}
\end{addmargin}
%-----------------------------------------------------------------------------
% Titel deiner Arbeit
%-----------------------------------------------------------------------------
\addvspace{2.6cm}
%TODO Formaenderung des Titels
\begin{addmargin}[\marginCoverPage]{0cm}
\begin{center}
{\fontfamily{phv}\selectfont\huge Bachelorarbeit} 
\end{center}

{\Large \par}
\addvspace{1.5cm}
\begin{center}
%TODO Wenn du ein "Sperrvermerk" für deine Arbeit brauchst, entferne den Kommentar in der naechsten Zeile
%\textbf{{\color{ercred}\fontfamily{phv}\selectfont{\huge Sperrvermerk}}}\\

%TODO  Der Englische Titel deiner Arbeit wird nach dem "\huge" in der naechsten Zeile geaendert
\textbf{\fontfamily{phv}\selectfont{\huge Design einer Abschlussarbeit}}
\end{center}

%\addvspace{1.5cm}
%\begin{center}
%TODO Falls der deutsche Name noch noetig ist, kommt er in die naechste Zeile
%{\fontfamily{phv}\selectfont Design of a Master's thesis}
%\end{center}

%-----------------------------------------------------------------------------
% Namen, Daten
%-----------------------------------------------------------------------------
\vfill
\begin{center}
\begingroup
\fontfamily{phv}\selectfont
%TODO Bitte aendere das Datum deiner Abgabe
Aachen, Mai 2019 \\
\addvspace{0.5cm}
%TODO Bitte aendere deinen Namen und Matrikelnummer in den naechsten beiden Zeilen
\textbf{Jonas Baumgärtner} \\
Matrikelnummer: 330575 \\
\addvspace{0.5cm}
%TODO Deine Betreuer. Achte auf die richtige Schreibweise des Titels.
Betreuer:\\
Sebastian Remy, M. Sc. \\
Univ.-Prof. Dr.-Ing. Dirk Müller \\
\addvspace{0.5cm}
%TODO  Eingabe der richtigen Anschrift des Instituts.
Diese Arbeit wurde vorgelegt am:\\
E.ON Energy Research Center | ERC \\
Institute for Energy Efficient Buildings and Indoor Climate | EBC\\
(Lehrstuhl für Gebäude- und Raumklimatechnik)\\
Mathieustraße 10, 52074 Aachen\\
\endgroup
\end{center}
\end{addmargin}
\end{titlepage}



%-----------------------------------------------------------------------------
% Setup von speziellen Seiten(Abbildungsverzeichnis, Nomenklatur etc.
%-----------------------------------------------------------------------------
\cleardoublepage
\pagenumbering{Roman} % Romanische Nummerierung
%-----------------------------------------------------------------------------
% Englische und Deutsche Kurzfassung. Beide notwendig
%-----------------------------------------------------------------------------
\clearpage
\refstepcounter{Hilfszaehler}
\renewcommand\abstractname{Kurzfassung}
\begin{abstract}
Überall dieselbe alte Leier. Das Layout ist fertig, der Text lässt auf sich warten. Damit das Layout nun nicht nackt im Raume steht und sich klein und leer vorkommt, springe ich ein: der Blindtext. Genau zu diesem Zwecke erschaffen, immer im Schatten meines großen Bruders »Lorem Ipsum«, freue ich mich jedes Mal, wenn Sie ein paar Zeilen lesen. Denn esse est percipi - Sein ist wahrgenommen werden.

Und weil Sie nun schon die Güte haben, mich ein paar weitere Sätze lang zu begleiten, möchte ich diese Gelegenheit nutzen, Ihnen nicht nur als Lückenfüller zu dienen, sondern auf etwas hinzuweisen, das es ebenso verdient wahrgenommen zu werden: Webstandards nämlich. Sehen Sie, Webstandards sind das Regelwerk, auf dem Webseiten aufbauen. So gibt es Regeln für HTML, CSS, JavaScript oder auch XML; Worte, die Sie vielleicht schon einmal von Ihrem Entwickler gehört haben. Diese Standards sorgen dafür, dass alle Beteiligten aus einer Webseite den größten Nutzen ziehen.

Im Gegensatz zu früheren Webseiten müssen wir zum Beispiel nicht mehr zwei verschiedene Webseiten für den Internet Explorer und einen anderen Browser programmieren. Es reicht eine Seite, die - richtig angelegt - sowohl auf verschiedenen Browsern im Netz funktioniert, aber ebenso gut für den Ausdruck oder die Darstellung auf einem Handy geeignet ist. Wohlgemerkt: Eine Seite für alle Formate. Was für eine Erleichterung. Standards sparen Zeit bei den Entwicklungskosten und sorgen dafür, dass sich Webseiten später leichter pflegen lassen. Natürlich nur dann, wenn sich alle an diese Standards halten. Das gilt für Browser wie Firefox, Opera, Safari und den Internet Explorer ebenso wie für die Darstellung in Handys. Und was können Sie für Standards tun? Fordern Sie von Ihren Designern und Programmieren einfach standardkonforme Webseiten. Ihr Budget wird es Ihnen auf Dauer danken. Ebenso möchte ich Ihnen dafür danken, dass Sie mich bin zum Ende gelesen

This abstract has 300 words

\end{abstract}
\renewcommand{\abstractname}{Abstract}
\begin{abstract}

 
\end{abstract}

%-----------------------------------------------------------------------------
% Inhaltsverzeichnis
%-----------------------------------------------------------------------------
\pagestyle{plain.scrheadings}
\setcounter{secnumdepth}{2}
\setcounter{tocdepth}{2}
\tableofcontents

%-----------------------------------------------------------------------------
% Nomenklatur
%-----------------------------------------------------------------------------
\clearpage
\refstepcounter{Hilfszaehler}
\addcontentsline{toc}{chapter}{Nomenklatur}
\chapter*{Nomenklatur}
\begin{onehalfspacing}
\begin{longtable}[h]{p{0.15\textwidth} p{0.65\textwidth} p{0.1\textwidth}}
		\caption*{\textbf{Formelzeichen und Einheiten}} \\
		\\
		\textbf{Symbol} & \textbf{Bedeutung} & \textbf{Einheit} \\ %\hline 
		\endhead
		\\
		\multicolumn{3}{c}{Fortsetzung auf der nächsten Seite} \\
		\endfoot
		\multicolumn{3}{c}{ } \\
		\endlastfoot
		
		$A$ & Fläche & \squaremetre\\
		$c_{p}$&spezifische Wärmekapazität bei konstantem Druck&\joule\per(\kilogram\usk\kelvin)\\
			$C$&Wärmekapazität&\watt\per\kilogram\\
		$H $ & Enthalpie & \joule\\		
		$\dot{H}$ & Enthalpiestrom & \joule\per\second\\
		$E$ & Exergie & \joule\\
		$e$ & spezifische Exergie & \joule\per\kilogram\\
		$\dot{m}$ & Massenstrom & \kilogram\per\second\\
		$p$ & Druck & \pascal\\
		$\dot{Q}$ & Wärmestrom & \watt\\
		$R$ & spezifische Gaskonstante & \joule\per(\kilogram\usk\kelvin)\\
		$S$ & Entropie & \joule\per\kelvin\\
		$\dot{S}$ & Entropiestrom & \watt\per\kelvin\\
		$T$ & Temperatur & \kelvin\\
		$t$ & Zeit & \second\\
		$U$ & innere Energie & \joule\\
		$U_{T}$ & Wärmedurchgangskoeffizient & \watt\per(\kilogram\usk\kelvin)\\
		$h$ & Wärmeübergangskoeffizient & \watt\per(\squaremetre\usk\kelvin)\\		
		$V$ & Volumen & \cubic\meter\\
		$\dot{V}$&Volumenstrom&\cubic\meter\per\second\\
		$\dot{W}$ & Leistung & \watt\\
		$Y$ & Wasserbeladung der Luft & \gram\per\kilogram\\
		
		
\end{longtable}

\begin{longtable}[h]{p{0.15\textwidth} p{0.65\textwidth} p{0.1\textwidth}}
		\caption*{\textbf{Griechische Formelzeichen}} \\
		\\
		\textbf{Symbol} & \textbf{Bedeutung} & \textbf{Einheit} \\ %\hline 
		\endhead
		\\
		\multicolumn{3}{c}{Fortsetzung auf der nächsten Seite} \\
		\endfoot
		\multicolumn{3}{c}{ } \\
		\endlastfoot
		
		$\eta_{C}$ & Carnot-Wirkungsgrad & ---\\
		$\kappa_{\mathrm{E}}$ & exergetische Aufwandszahl der Wärmeerzeugung & ---\\
		$\kappa_{\mathrm{T}}$ & exergetische Aufwandszahl des Wärmetransfers & ---\\
		$\Phi$ & thermische Leistung & \watt\\
		$\varrho$& Massendichte&\kilogrampercubicmetre\\
			$\sigma$&Temperaturspreizung&\kelvin\\
		$\vartheta $ & Temperatur  & \degreecelsius\\
		$\Delta\vartheta $ & Temperaturdifferenz  &\kelvin\\
		
\end{longtable}

\begin{longtable}[h]{p{0.15\textwidth} p{0.75\textwidth}}
		\caption*{\textbf{Indizes und Abkürzungen}} \\
		\\
		\textbf{Symbol} & \textbf{Bedeutung} \\ %\hline 
		\endhead
		\\
		\multicolumn{2}{c}{Fortsetzung auf der nächsten Seite} \\
		\endfoot
		\multicolumn{2}{c}{ } \\
		\endlastfoot
		
		0 & Referenzzustand (\emph{ambient dead state})\\
		A & Außen/Umgebung\\ 		
		CH & chemisch\\
		CV & Kontrollvolumen (\emph{control volume})\\
		DSC & Dynamische Differenzkalorimetrie (\emph{differential scanning calorimetry}) \\
		e & über die Systemgrenze (\emph{external})\\
		F & Volumenstrom\\	
		FW & Fassadenwärmeübertrager\\
		gen & erzeugt (\emph{generated})\\
		In & Eingang (\emph{input})\\
		KN & kinetisch\\
		KRM & Kapillarrohrmatte\\
		LabVIEW & Programmiersprache und Entwicklungsumgebung für die Messdatenerfassung der Firma National Instruments\\
		L & Luft\\
		LWS & Latentwärmespeicher\\	
		m & Mittelwert\\
			Ob&Oberfläche\\
		PCM & Latentwärmespeichermaterial (\emph{phase change material})\\
		PH & physikalisch\\
		PT & potentiell\\
		Q & auf einen Wärmestrom bezogen\\
		R & Rücklauf\\
		Reg& Speicherregeneration\\
		T & Temperatur\\
		$\Delta$ t & Zeitschritt der Länge $\Delta$ t\\
		t & technisch\\
		V & Vorlauf\\
		V & Verlust (Exergieanalyse)\\
		W&Wärmeträgermedium\\
		
\end{longtable}
\end{onehalfspacing}


%-----------------------------------------------------------------------------
% Abbildungsverzeichnis
%-----------------------------------------------------------------------------
\clearpage
\refstepcounter{Hilfszaehler}
\addcontentsline{toc}{chapter}{Abbildungsverzeichnis}
\listoffigures

%-----------------------------------------------------------------------------
% Tabellenverzeichnis
%-----------------------------------------------------------------------------
\clearpage
\refstepcounter{Hilfszaehler}
\addcontentsline{toc}{chapter}{Tabellenverzeichnis}
\listoftables

%-----------------------------------------------------------------------------
% Vorwort, kommentiere die Zeile wenn nicht nötig
%-----------------------------------------------------------------------------
\clearpage
\refstepcounter{Hilfszaehler}
\addcontentsline{toc}{chapter}{Vorwort}
%\newpage\thispagestyle{empty}~
\chapter*{Vorwort}



%-----------------------------------------------------------------------------
% Setup Normale Seiten(Neue Nummerierung)
%-----------------------------------------------------------------------------
\cleardoublepage
\clearpage
\pagenumbering{arabic}
\pagestyle{scrheadings}%\ihead{\leftmark}\chead{}\ohead{\rightmark}
\automark[chapter]{chapter}
\automark*[section]{}

%-----------------------------------------------------------------------------
% Füge hier deine Kapitel hinzu
%-----------------------------------------------------------------------------
\chapter{Einleitung}

\section{Motivation}
\label{sec:Sektion 11}
das ist die einleitung
\section{Aufgabenstellung}
\label{sec:Sektion 12}
das ist die aufgabenstellung
%\chapter{Stand der Forschung}

\section{Systemische Bewertung}

\section{Stromnetz}

\subsection{Netztopologien}
\subsection{Betriebsmittel}

\chapter{xyz}

\section{Motivation}
\label{sec:Sektion 1}
Das ist ein Test, wie ich meine BA mit LaTex und GitHub schreiben kann.
Klappt es?
Man weiß es nicht...
\section{Aufgabenstellung}
\label{sec:Sektion 2}
jhnklknhh


\chapter{neuesKapitel}

\section{Motivation}
\label{sec:Sektion 1}
fjkklhkjlkll
\section{Aufgabenstellung}
\label{sec:Sektion 2}
jhnklknhh
\subsection{Gebäudepups}

%-----------------------------------------------------------------------------
% Literaturverzeichnis
%-----------------------------------------------------------------------------

\clearpage
\refstepcounter{Hilfszaehler}
\addcontentsline{toc}{chapter}{\bibname}
\protect\bibliographystyle{natdin} % Dieser Stil ist eine Empfehlung, aber nicht verprlichtend.
\bibliography{Literaturverzeichnis}

%-----------------------------------------------------------------------------
% Anhang
%-----------------------------------------------------------------------------
\pagestyle{empty}
\clearpage
\vfill{}\vspace{-1\parskip}
\begin{center}\textsf{\textbf{\huge Anhang}}\end{center}{\huge \par}
\vfill{}
\appendix
\newpage
\pagestyle{scrheadings}
\chapter{Tabellen}


\section{U-Werte nach TABULA}

\begin{table}[H]\centering
\begin{tabular}{l|cccc|cccc}
\toprule[1.5pt]
 & \multicolumn{4}{c}{SFH}& \multicolumn{4}{c}{MFH}  \\ 
\cmidrule[1.5pt]{2-9}
Baualtersklasse & Dach & Außenwand & Fenster & Boden & Dach & Außenwand & Fenster & Boden \\ \addlinespace[5pt]
\midrule[2pt]
vor 1918 & 1,3 & 1,7 & 2,8 & 0,88 & 1,3 & 2,2 & 2,7 & 0,88 \\
\midrule
1919 - 1948 & 1,4 & 1,7 & 2,8 & 0,77 & 1,4 & 1,7 & 3 & 0,77 \\
\midrule
1949 - 1957 & 1,4 & 1,4 & 2,8 & 0,78 & 1,08 & 1,2 & 3 & 1,33 \\
\midrule
1958 - 1968 & 0,8 & 1,2 & 2,8 & 1,08 & 0,51 & 1,2 & 3 & 1,08 \\
\midrule
1969 - 1978 & 0,5 & 1 & 2,8 & 0,77 & 0,51 & 1 & 3 & 0,77 \\
\midrule
1979 - 1983 & 0,5 & 0,8 & 4,3 & 0,65 & 0,43 & 0,8 & 3 & 0,65 \\
\midrule
1984 - 1994 & 0,4 & 0,5 & 3,2 & 0,51 & 0,36 & 0,6 & 3 & 0,51 \\
\midrule
1995 - 2001 & 0,35 & 0,3 & 1,9 & 0,4 & 0,32 & 0,4 & 1,9 & 0,4\\
\bottomrule[1.5pt] \addlinespace[10pt]
\end{tabular}
\caption{Wärmedurchgangskoeffizienten der Bauteile Dach, Außenwand, Fenster und Boden nach Gebäudeart und Baualtersklasse [in \(W/(m^2 \cdot K)\)]}
\label{tab: TabelleA1}
\end{table}


\chapter{Wichtiger Anhang}


%-----------------------------------------------------------------------------
% Eigenstaendigkeitserklaerung
%-----------------------------------------------------------------------------
\pagestyle{empty}
\cleardoublepage
\Large
\textbf{Eigenständigkeitserklärung}

\normalsize
Hiermit versichere ich, dass ich die vorliegende Arbeit selbstständig und ohne Benutzung anderer als der angegebenen Hilfsmittel angefertigt habe. Alle Stellen, die wörtlich oder sinngemäß übernommen sind, sind als solche kenntlich gemacht. Die Arbeit ist in gleicher oder ähnlicher Form noch nicht als Prüfungsarbeit eingereicht worden. Ich erkläre mich damit einverstanden, dass die vorliegende Arbeit in der Lehrstuhlbibliothek und Datenbank aufbewahrt und für den internen Gebrauch kopiert werden darf. \newline \newline


Aachen, den \today \newline \newline

%TODO Name ergänzen
DEIN NAME



\end{document}